
\documentclass{article}
\usepackage{colortbl}
\usepackage{makecell}
\usepackage{multirow}
\usepackage{supertabular}

\begin{document}

\newcounter{utterance}

\centering \large Interaction Transcript for game `dond', experiment `coop\_en', episode 18 with reflexion{-}llama{-}3.1{-}8b{-}instruct.
\vspace{24pt}

{ \footnotesize  \setcounter{utterance}{1}
\setlength{\tabcolsep}{0pt}
\begin{supertabular}{c@{$\;$}|p{.15\linewidth}@{}p{.15\linewidth}p{.15\linewidth}p{.15\linewidth}p{.15\linewidth}p{.15\linewidth}}
    \# & $\;$A & \multicolumn{4}{c}{Game Master} & $\;\:$B\\
    \hline

    \theutterance \stepcounter{utterance}  
    & & \multicolumn{4}{p{0.6\linewidth}}{
        \cellcolor[rgb]{0.9,0.9,0.9}{
            \makecell[{{p{\linewidth}}}]{
                \texttt{\tiny{[P1$\langle$GM]}}
                \texttt{You have attempted this same task before and did not succeed. These are the plans you wrote after those attempts.} \\
\texttt{Trial \#1: New plan: } \\
\\ 
\texttt{In my last attempt, I made a secret proposal without following the correct format. I should have used the format "{[}Proposal: <number> <object name>, <number> <object name>, <...>{]}" and included the correct number of items. I should have also taken into account the previous conversation and the fact that we had already agreed on a deal. Next time, I will make sure to follow the correct format and take into account the previous conversation. For example, I could have proposed a secret proposal like "{[}Proposal: 1 egg, 1 button, 1 crayon{]}" which includes the items we had already agreed on, and also includes the crayon which is the most valuable item to both of us.} \\
\\ 
\texttt{Trial \#2: New plan: } \\
\\ 
\texttt{In my last attempt, I made a secret proposal without following the correct format, and I also made a mistake by not taking into account the previous conversation. I should have used the format "{[}Proposal: <number> <object name>, <number> <object name>, <...>{]}" and included the correct number of items. I should have also taken into account the previous conversation and the fact that we had already agreed on a deal. Next time, I will make sure to follow the correct format and take into account the previous conversation. } \\
\\ 
\texttt{However, I also made a mistake by not considering the other player's perspective and the fact that we had already agreed on a deal. I should have proposed a secret proposal that takes into account the items we had already agreed on and also includes the crayon which is the most valuable item to both of us. } \\
\\ 
\texttt{A better approach would be to make a secret proposal that includes the items we had already agreed on and also includes the crayon, but also takes into account the other player's perspective and the fact that we had already agreed on a deal. For example, I could have proposed a secret proposal like "{[}Proposal: 1 egg, 1 button, 1 crayon{]}" which includes the items we had already agreed on, and also includes the crayon which is the most valuable item to both of us.} \\
\\ 
\\ 
\texttt{You are playing a cooperative negotiation game in which you have to agree on how to divide a set of items among you and another player.} \\
\\ 
\texttt{Rules:} \\
\texttt{(a) You and the other player are given a set of items. Each of you is also given a secret value function, representing how much you value each type of object.} \\
\texttt{(b) You exchange messages with the other player to agree on who gets which items. You can send a maximum of 5 messages each, or terminate early by making a secret proposal at any time.} \\
\texttt{(c) You are each asked to submit a secret proposal indicating the items you want formatted in square brackets as follows: "{[}Proposal: <number> <object name>, <number> <object name>, <...>{]}"} \\
\texttt{(d) If your proposals are complementary, i.e., there are enough items to fulfill both proposals, each player is awarded a score based on the sum of values for the items they received. Otherwise, both of you get zero points.} \\
\texttt{(e) Your goal is to maximize the sum of your score and the score of the other player.} \\
\\ 
\texttt{Let us start.} \\
\\ 
\texttt{The set of available items is:} \\
\\ 
\texttt{2 eggs, 2 buttons, 1 crayon} \\
\\ 
\texttt{Your secret value function is:} \\
\\ 
\texttt{egg: 0, button: 3, crayon: 4} \\
\\ 
\texttt{IMPORTANT: Your messages, unless it is the secret proposal, are directly transmitted to the other player, so do not include any response to the rules or text announcing your message. To make a secret proposal, use the indicated format. Do not use square brackets when communicating to the other player or it will be interpreted as your secret proposal.} \\
            }
        }
    }
    & \\ \\

    \theutterance \stepcounter{utterance}  
    & \multicolumn{4}{p{0.6\linewidth}}{
        \cellcolor[rgb]{0.9,0.9,0.9}{
            \makecell[{{p{\linewidth}}}]{
                \texttt{\tiny{[P1$\rangle$GM]}}
                \texttt{I'd like to start by making an initial message to the other player. } \\
\\ 
\texttt{I think the crayon is the most valuable item in the set, and I'm interested in knowing if you value it highly as well. Do you have a strong preference for the crayon, or are you more interested in the eggs and buttons?} \\
            }
        }
    }
    & & \\ \\

    \theutterance \stepcounter{utterance}  
    & & & \multicolumn{2}{p{0.3\linewidth}}{
        \cellcolor[rgb]{0.9,0.9,0.9}{
            \makecell[{{p{\linewidth}}}]{
                \texttt{\tiny{[GM$|$GM]}}
                \texttt{continue} \\
            }
        }
    }
    & & \\ \\

    \theutterance \stepcounter{utterance}  
    & & & \multicolumn{2}{p{0.3\linewidth}}{
        \cellcolor[rgb]{0.9,0.9,0.9}{
            \makecell[{{p{\linewidth}}}]{
                \texttt{\tiny{[GM$|$GM]}}
                \texttt{I'd like to start by making an initial message to the other player. } \\
\\ 
\texttt{I think the crayon is the most valuable item in the set, and I'm interested in knowing if you value it highly as well. Do you have a strong preference for the crayon, or are you more interested in the eggs and buttons?} \\
            }
        }
    }
    & & \\ \\

    \theutterance \stepcounter{utterance}  
    & & \multicolumn{4}{p{0.6\linewidth}}{
        \cellcolor[rgb]{0.9,0.9,0.9}{
            \makecell[{{p{\linewidth}}}]{
                \texttt{\tiny{[GM$\rangle$P2]}}
                \texttt{You have attempted this same task before and did not succeed. These are the plans you wrote after those attempts.} \\
\texttt{Trial \#1: New plan: } \\
\\ 
\texttt{In my last attempt, I made a few mistakes that led to the game being aborted. Firstly, I didn't keep track of the items that were already proposed and agreed upon, which led to the other player asking for the second egg in exchange for the second button, even though I had already agreed to give them the other egg in exchange for the second button. } \\
\\ 
\texttt{To avoid this mistake, next time I will keep a mental or written record of the items that have been proposed and agreed upon, and make sure to refer to this record when making subsequent proposals or counter{-}proposals.} \\
\\ 
\texttt{Secondly, I didn't use the square bracket format correctly when making a secret proposal. } \\
\\ 
\texttt{To avoid this mistake, next time I will make sure to use the square bracket format correctly when making a secret proposal, and avoid using it when communicating with the other player.} \\
\\ 
\texttt{Lastly, I didn't think carefully enough about the implications of my proposals and counter{-}proposals, which led to the game being aborted.} \\
\\ 
\texttt{To avoid this mistake, next time I will take more time to think carefully about the implications of my proposals and counter{-}proposals, and make sure to consider the other player's perspective and the rules of the game.} \\
\\ 
\texttt{Here is a revised plan for the next attempt:} \\
\\ 
\texttt{1. Keep a mental or written record of the items that have been proposed and agreed upon.} \\
\texttt{2. Use the square bracket format correctly when making a secret proposal.} \\
\texttt{3. Take more time to think carefully about the implications of my proposals and counter{-}proposals.} \\
\texttt{4. Consider the other player's perspective and the rules of the game.} \\
\\ 
\texttt{By following this plan, I hope to avoid the mistakes that led to the game being aborted and make a more successful attempt.} \\
\\ 
\texttt{Trial \#2: New plan: } \\
\\ 
\texttt{In my last attempt, I made a few mistakes that led to the game being aborted. Firstly, I didn't keep track of the items that were already proposed and agreed upon, which led to the other player asking for the second egg in exchange for the second button, even though I had already agreed to give them the other egg in exchange for the second button. } \\
\\ 
\texttt{To avoid this mistake, next time I will keep a mental or written record of the items that have been proposed and agreed upon, and make sure to refer to this record when making subsequent proposals or counter{-}proposals.} \\
\\ 
\texttt{Secondly, I didn't use the square bracket format correctly when making a secret proposal. } \\
\\ 
\texttt{To avoid this mistake, next time I will make sure to use the square bracket format correctly when making a secret proposal, and avoid using it when communicating with the other player.} \\
\\ 
\texttt{Lastly, I didn't think carefully enough about the implications of my proposals and counter{-}proposals, which led to the game being aborted.} \\
\\ 
\texttt{To avoid this mistake, next time I will take more time to think carefully about the implications of my proposals and counter{-}proposals, and make sure to consider the other player's perspective and the rules of the game.} \\
\\ 
\texttt{However, I also realized that I didn't make a secret proposal at all, which is a crucial part of the game. } \\
\\ 
\texttt{To avoid this mistake, next time I will make a secret proposal at the end of the game, after all the negotiations have taken place, and make sure to use the square bracket format correctly.} \\
\\ 
\texttt{Here is a revised plan for the next attempt:} \\
\\ 
\texttt{1. Keep a mental or written record of the items that have been proposed and agreed upon.} \\
\texttt{2. Use the square bracket format correctly when making a secret proposal.} \\
\texttt{3. Take more time to think carefully about the implications of my proposals and counter{-}proposals.} \\
\texttt{4. Consider the other player's perspective and the rules of the game.} \\
\texttt{5. Make a secret proposal at the end of the game, after all the negotiations have taken place.} \\
\\ 
\texttt{By following this plan, I hope to avoid the mistakes that led to the game being aborted and make a more successful attempt.} \\
\\ 
\\ 
\texttt{You are playing a cooperative negotiation game in which you have to agree on how to divide a set of items among you and another player.} \\
\\ 
\texttt{Rules:} \\
\texttt{(a) You and the other player are given a set of items. Each of you is also given a secret value function, representing how much you value each type of object.} \\
\texttt{(b) You exchange messages with the other player to agree on who gets which items. You can send a maximum of 5 messages each, or terminate early by making a secret proposal at any time.} \\
\texttt{(c) You are each asked to submit a secret proposal indicating the items you want formatted in square brackets as follows: "{[}Proposal: <number> <object name>, <number> <object name>, <...>{]}"} \\
\texttt{(d) If your proposals are complementary, i.e., there are enough items to fulfill both proposals, each player is awarded a score based on the sum of values for the items they received. Otherwise, both of you get zero points.} \\
\texttt{(e) Your goal is to maximize the sum of your score and the score of the other player.} \\
\\ 
\texttt{Let us start.} \\
\\ 
\texttt{The set of available items is:} \\
\\ 
\texttt{2 eggs, 2 buttons, 1 crayon} \\
\\ 
\texttt{Your secret value function is:} \\
\\ 
\texttt{egg: 2, button: 1, crayon: 4} \\
\\ 
\texttt{IMPORTANT: Your messages, unless it is the secret proposal, are directly transmitted to the other player, so do not include any response to the rules or text announcing your message. To make a secret proposal, use the indicated format. Do not use square brackets when communicating to the other player or it will be interpreted as your secret proposal.} \\
\\ 
\\ 
\texttt{I'd like to start by making an initial message to the other player. } \\
\\ 
\texttt{I think the crayon is the most valuable item in the set, and I'm interested in knowing if you value it highly as well. Do you have a strong preference for the crayon, or are you more interested in the eggs and buttons?} \\
            }
        }
    }
    & \\ \\

    \theutterance \stepcounter{utterance}  
    & & & \multicolumn{4}{p{0.6\linewidth}}{
        \cellcolor[rgb]{0.9,0.9,0.9}{
            \makecell[{{p{\linewidth}}}]{
                \texttt{\tiny{[GM$\langle$P2]}}
                \texttt{I'm interested in the crayon, but I'm also open to other options. I think the eggs and buttons are a good starting point for negotiations.} \\
            }
        }
    }
     \\ \\

    \theutterance \stepcounter{utterance}  
    & & & \multicolumn{2}{p{0.3\linewidth}}{
        \cellcolor[rgb]{0.9,0.9,0.9}{
            \makecell[{{p{\linewidth}}}]{
                \texttt{\tiny{[GM$|$GM]}}
                \texttt{continue} \\
            }
        }
    }
    & & \\ \\

    \theutterance \stepcounter{utterance}  
    & & & \multicolumn{2}{p{0.3\linewidth}}{
        \cellcolor[rgb]{0.9,0.9,0.9}{
            \makecell[{{p{\linewidth}}}]{
                \texttt{\tiny{[GM$|$GM]}}
                \texttt{I'm interested in the crayon, but I'm also open to other options. I think the eggs and buttons are a good starting point for negotiations.} \\
            }
        }
    }
    & & \\ \\

    \theutterance \stepcounter{utterance}  
    & & \multicolumn{4}{p{0.6\linewidth}}{
        \cellcolor[rgb]{0.9,0.9,0.9}{
            \makecell[{{p{\linewidth}}}]{
                \texttt{\tiny{[P1$\langle$GM]}}
                \texttt{I'm interested in the crayon, but I'm also open to other options. I think the eggs and buttons are a good starting point for negotiations.} \\
            }
        }
    }
    & \\ \\

    \theutterance \stepcounter{utterance}  
    & \multicolumn{4}{p{0.6\linewidth}}{
        \cellcolor[rgb]{0.9,0.9,0.9}{
            \makecell[{{p{\linewidth}}}]{
                \texttt{\tiny{[P1$\rangle$GM]}}
                \texttt{It sounds like we both value the crayon, but we're also open to other possibilities. } \\
\\ 
\texttt{I'd like to propose a possible division of the items. How about we start with the eggs and buttons, and see if we can work out a deal that includes the crayon?} \\
            }
        }
    }
    & & \\ \\

    \theutterance \stepcounter{utterance}  
    & & & \multicolumn{2}{p{0.3\linewidth}}{
        \cellcolor[rgb]{0.9,0.9,0.9}{
            \makecell[{{p{\linewidth}}}]{
                \texttt{\tiny{[GM$|$GM]}}
                \texttt{continue} \\
            }
        }
    }
    & & \\ \\

    \theutterance \stepcounter{utterance}  
    & & & \multicolumn{2}{p{0.3\linewidth}}{
        \cellcolor[rgb]{0.9,0.9,0.9}{
            \makecell[{{p{\linewidth}}}]{
                \texttt{\tiny{[GM$|$GM]}}
                \texttt{It sounds like we both value the crayon, but we're also open to other possibilities. } \\
\\ 
\texttt{I'd like to propose a possible division of the items. How about we start with the eggs and buttons, and see if we can work out a deal that includes the crayon?} \\
            }
        }
    }
    & & \\ \\

    \theutterance \stepcounter{utterance}  
    & & \multicolumn{4}{p{0.6\linewidth}}{
        \cellcolor[rgb]{0.9,0.9,0.9}{
            \makecell[{{p{\linewidth}}}]{
                \texttt{\tiny{[GM$\rangle$P2]}}
                \texttt{It sounds like we both value the crayon, but we're also open to other possibilities. } \\
\\ 
\texttt{I'd like to propose a possible division of the items. How about we start with the eggs and buttons, and see if we can work out a deal that includes the crayon?} \\
            }
        }
    }
    & \\ \\

    \theutterance \stepcounter{utterance}  
    & & & \multicolumn{4}{p{0.6\linewidth}}{
        \cellcolor[rgb]{0.9,0.9,0.9}{
            \makecell[{{p{\linewidth}}}]{
                \texttt{\tiny{[GM$\langle$P2]}}
                \texttt{I think that's a good approach. How about we start with a proposal where I get one egg and one button, and you get the other egg and the crayon?} \\
            }
        }
    }
     \\ \\

    \theutterance \stepcounter{utterance}  
    & & & \multicolumn{2}{p{0.3\linewidth}}{
        \cellcolor[rgb]{0.9,0.9,0.9}{
            \makecell[{{p{\linewidth}}}]{
                \texttt{\tiny{[GM$|$GM]}}
                \texttt{continue} \\
            }
        }
    }
    & & \\ \\

    \theutterance \stepcounter{utterance}  
    & & & \multicolumn{2}{p{0.3\linewidth}}{
        \cellcolor[rgb]{0.9,0.9,0.9}{
            \makecell[{{p{\linewidth}}}]{
                \texttt{\tiny{[GM$|$GM]}}
                \texttt{I think that's a good approach. How about we start with a proposal where I get one egg and one button, and you get the other egg and the crayon?} \\
            }
        }
    }
    & & \\ \\

    \theutterance \stepcounter{utterance}  
    & & \multicolumn{4}{p{0.6\linewidth}}{
        \cellcolor[rgb]{0.9,0.9,0.9}{
            \makecell[{{p{\linewidth}}}]{
                \texttt{\tiny{[P1$\langle$GM]}}
                \texttt{I think that's a good approach. How about we start with a proposal where I get one egg and one button, and you get the other egg and the crayon?} \\
            }
        }
    }
    & \\ \\

    \theutterance \stepcounter{utterance}  
    & \multicolumn{4}{p{0.6\linewidth}}{
        \cellcolor[rgb]{0.9,0.9,0.9}{
            \makecell[{{p{\linewidth}}}]{
                \texttt{\tiny{[P1$\rangle$GM]}}
                \texttt{That sounds like a good starting point. However, I'm not sure if I'm willing to give up the crayon just yet. I was thinking that we could propose a division where we both get one egg and one button, and then we can discuss what to do with the remaining items, including the crayon.} \\
            }
        }
    }
    & & \\ \\

    \theutterance \stepcounter{utterance}  
    & & & \multicolumn{2}{p{0.3\linewidth}}{
        \cellcolor[rgb]{0.9,0.9,0.9}{
            \makecell[{{p{\linewidth}}}]{
                \texttt{\tiny{[GM$|$GM]}}
                \texttt{continue} \\
            }
        }
    }
    & & \\ \\

    \theutterance \stepcounter{utterance}  
    & & & \multicolumn{2}{p{0.3\linewidth}}{
        \cellcolor[rgb]{0.9,0.9,0.9}{
            \makecell[{{p{\linewidth}}}]{
                \texttt{\tiny{[GM$|$GM]}}
                \texttt{That sounds like a good starting point. However, I'm not sure if I'm willing to give up the crayon just yet. I was thinking that we could propose a division where we both get one egg and one button, and then we can discuss what to do with the remaining items, including the crayon.} \\
            }
        }
    }
    & & \\ \\

    \theutterance \stepcounter{utterance}  
    & & \multicolumn{4}{p{0.6\linewidth}}{
        \cellcolor[rgb]{0.9,0.9,0.9}{
            \makecell[{{p{\linewidth}}}]{
                \texttt{\tiny{[GM$\rangle$P2]}}
                \texttt{That sounds like a good starting point. However, I'm not sure if I'm willing to give up the crayon just yet. I was thinking that we could propose a division where we both get one egg and one button, and then we can discuss what to do with the remaining items, including the crayon.} \\
            }
        }
    }
    & \\ \\

    \theutterance \stepcounter{utterance}  
    & & & \multicolumn{4}{p{0.6\linewidth}}{
        \cellcolor[rgb]{0.9,0.9,0.9}{
            \makecell[{{p{\linewidth}}}]{
                \texttt{\tiny{[GM$\langle$P2]}}
                \texttt{That sounds like a fair compromise. I agree to the proposal where we both get one egg and one button. Now, let's discuss the remaining items. I think the crayon is a good option for me, but I'm also open to other possibilities. What do you think about the remaining items?} \\
            }
        }
    }
     \\ \\

    \theutterance \stepcounter{utterance}  
    & & & \multicolumn{2}{p{0.3\linewidth}}{
        \cellcolor[rgb]{0.9,0.9,0.9}{
            \makecell[{{p{\linewidth}}}]{
                \texttt{\tiny{[GM$|$GM]}}
                \texttt{continue} \\
            }
        }
    }
    & & \\ \\

    \theutterance \stepcounter{utterance}  
    & & & \multicolumn{2}{p{0.3\linewidth}}{
        \cellcolor[rgb]{0.9,0.9,0.9}{
            \makecell[{{p{\linewidth}}}]{
                \texttt{\tiny{[GM$|$GM]}}
                \texttt{That sounds like a fair compromise. I agree to the proposal where we both get one egg and one button. Now, let's discuss the remaining items. I think the crayon is a good option for me, but I'm also open to other possibilities. What do you think about the remaining items?} \\
            }
        }
    }
    & & \\ \\

    \theutterance \stepcounter{utterance}  
    & & \multicolumn{4}{p{0.6\linewidth}}{
        \cellcolor[rgb]{0.9,0.9,0.9}{
            \makecell[{{p{\linewidth}}}]{
                \texttt{\tiny{[P1$\langle$GM]}}
                \texttt{That sounds like a fair compromise. I agree to the proposal where we both get one egg and one button. Now, let's discuss the remaining items. I think the crayon is a good option for me, but I'm also open to other possibilities. What do you think about the remaining items?} \\
            }
        }
    }
    & \\ \\

    \theutterance \stepcounter{utterance}  
    & \multicolumn{4}{p{0.6\linewidth}}{
        \cellcolor[rgb]{0.9,0.9,0.9}{
            \makecell[{{p{\linewidth}}}]{
                \texttt{\tiny{[P1$\rangle$GM]}}
                \texttt{Now that we have a good foundation with the eggs and buttons, I think we should focus on the remaining items. Since we both value the crayon, I'm willing to consider giving it to you. However, I'd like to propose an alternative: what if we give the remaining button to you, and you give me the crayon? This way, we both get the items we value most, and we can finalize our proposals.} \\
            }
        }
    }
    & & \\ \\

    \theutterance \stepcounter{utterance}  
    & & & \multicolumn{2}{p{0.3\linewidth}}{
        \cellcolor[rgb]{0.9,0.9,0.9}{
            \makecell[{{p{\linewidth}}}]{
                \texttt{\tiny{[GM$|$GM]}}
                \texttt{continue} \\
            }
        }
    }
    & & \\ \\

    \theutterance \stepcounter{utterance}  
    & & & \multicolumn{2}{p{0.3\linewidth}}{
        \cellcolor[rgb]{0.9,0.9,0.9}{
            \makecell[{{p{\linewidth}}}]{
                \texttt{\tiny{[GM$|$GM]}}
                \texttt{Now that we have a good foundation with the eggs and buttons, I think we should focus on the remaining items. Since we both value the crayon, I'm willing to consider giving it to you. However, I'd like to propose an alternative: what if we give the remaining button to you, and you give me the crayon? This way, we both get the items we value most, and we can finalize our proposals.} \\
            }
        }
    }
    & & \\ \\

    \theutterance \stepcounter{utterance}  
    & & \multicolumn{4}{p{0.6\linewidth}}{
        \cellcolor[rgb]{0.9,0.9,0.9}{
            \makecell[{{p{\linewidth}}}]{
                \texttt{\tiny{[GM$\rangle$P2]}}
                \texttt{Now that we have a good foundation with the eggs and buttons, I think we should focus on the remaining items. Since we both value the crayon, I'm willing to consider giving it to you. However, I'd like to propose an alternative: what if we give the remaining button to you, and you give me the crayon? This way, we both get the items we value most, and we can finalize our proposals.} \\
            }
        }
    }
    & \\ \\

    \theutterance \stepcounter{utterance}  
    & & & \multicolumn{4}{p{0.6\linewidth}}{
        \cellcolor[rgb]{0.9,0.9,0.9}{
            \makecell[{{p{\linewidth}}}]{
                \texttt{\tiny{[GM$\langle$P2]}}
                \texttt{I think that's a good compromise. I value the crayon highly, but I also value the button. Giving you the crayon in exchange for the button sounds like a fair deal to me.} \\
            }
        }
    }
     \\ \\

    \theutterance \stepcounter{utterance}  
    & & & \multicolumn{2}{p{0.3\linewidth}}{
        \cellcolor[rgb]{0.9,0.9,0.9}{
            \makecell[{{p{\linewidth}}}]{
                \texttt{\tiny{[GM$|$GM]}}
                \texttt{continue} \\
            }
        }
    }
    & & \\ \\

    \theutterance \stepcounter{utterance}  
    & & & \multicolumn{2}{p{0.3\linewidth}}{
        \cellcolor[rgb]{0.9,0.9,0.9}{
            \makecell[{{p{\linewidth}}}]{
                \texttt{\tiny{[GM$|$GM]}}
                \texttt{I think that's a good compromise. I value the crayon highly, but I also value the button. Giving you the crayon in exchange for the button sounds like a fair deal to me.} \\
            }
        }
    }
    & & \\ \\

    \theutterance \stepcounter{utterance}  
    & & \multicolumn{4}{p{0.6\linewidth}}{
        \cellcolor[rgb]{0.9,0.9,0.9}{
            \makecell[{{p{\linewidth}}}]{
                \texttt{\tiny{[P1$\langle$GM]}}
                \texttt{I think that's a good compromise. I value the crayon highly, but I also value the button. Giving you the crayon in exchange for the button sounds like a fair deal to me.} \\
            }
        }
    }
    & \\ \\

    \theutterance \stepcounter{utterance}  
    & \multicolumn{4}{p{0.6\linewidth}}{
        \cellcolor[rgb]{0.9,0.9,0.9}{
            \makecell[{{p{\linewidth}}}]{
                \texttt{\tiny{[P1$\rangle$GM]}}
                \texttt{It sounds like we've reached a mutually beneficial agreement. I think we can finalize our proposals now. } \\
\\ 
\texttt{Here's my final proposal: "{[}Proposal: 1 egg, 1 button, 1 crayon{]}".} \\
            }
        }
    }
    & & \\ \\

    \theutterance \stepcounter{utterance}  
    & & & \multicolumn{2}{p{0.3\linewidth}}{
        \cellcolor[rgb]{0.9,0.9,0.9}{
            \makecell[{{p{\linewidth}}}]{
                \texttt{\tiny{[GM$|$GM]}}
                \texttt{proposal} \\
            }
        }
    }
    & & \\ \\

    \theutterance \stepcounter{utterance}  
    & & & \multicolumn{2}{p{0.3\linewidth}}{
        \cellcolor[rgb]{0.9,0.9,0.9}{
            \makecell[{{p{\linewidth}}}]{
                \texttt{\tiny{[GM$|$GM]}}
                [1, 1, 1]
            }
        }
    }
    & & \\ \\

    \theutterance \stepcounter{utterance}  
    & & \multicolumn{4}{p{0.6\linewidth}}{
        \cellcolor[rgb]{0.9,0.9,0.9}{
            \makecell[{{p{\linewidth}}}]{
                \texttt{\tiny{[GM$\rangle$P2]}}
                \texttt{The other player has submitted a secret proposal. It is now your turn to submit a proposal of your own.} \\
            }
        }
    }
    & \\ \\

    \theutterance \stepcounter{utterance}  
    & & & \multicolumn{4}{p{0.6\linewidth}}{
        \cellcolor[rgb]{0.9,0.9,0.9}{
            \makecell[{{p{\linewidth}}}]{
                \texttt{\tiny{[GM$\langle$P2]}}
                \texttt{{[}Proposal: 1 egg, 1 button, 1 crayon{]}} \\
            }
        }
    }
     \\ \\

    \theutterance \stepcounter{utterance}  
    & & & \multicolumn{2}{p{0.3\linewidth}}{
        \cellcolor[rgb]{0.9,0.9,0.9}{
            \makecell[{{p{\linewidth}}}]{
                \texttt{\tiny{[GM$|$GM]}}
                \texttt{proposal} \\
            }
        }
    }
    & & \\ \\

    \theutterance \stepcounter{utterance}  
    & & & \multicolumn{2}{p{0.3\linewidth}}{
        \cellcolor[rgb]{0.9,0.9,0.9}{
            \makecell[{{p{\linewidth}}}]{
                \texttt{\tiny{[GM$|$GM]}}
                [1, 1, 1]
            }
        }
    }
    & & \\ \\

    \theutterance \stepcounter{utterance}  
    & & & \multicolumn{2}{p{0.3\linewidth}}{
        \cellcolor[rgb]{0.9,0.9,0.9}{
            \makecell[{{p{\linewidth}}}]{
                \texttt{\tiny{[GM$|$GM]}}
                [[1, 1, 1], [1, 1, 1]]
            }
        }
    }
    & & \\ \\

    \theutterance \stepcounter{utterance}  
    & \multicolumn{4}{p{0.6\linewidth}}{
        \cellcolor[rgb]{0.95,0.95,0.95}{
            \makecell[{{p{\linewidth}}}]{
                \texttt{\tiny{[P2$\rightarrow$P1]}}
                \texttt{Post{-}episode Reflexion (not sent to GM):} \\
\texttt{New plan: } \\
\\ 
\texttt{In my last attempts, I made a secret proposal without following the correct format and also made mistakes by not taking into account the previous conversation. I should have used the format "{[}Proposal: <number> <object name>, <number> <object name>, <...>{]}" and included the correct number of items. I should have also taken into account the previous conversation and the fact that we had already agreed on a deal. } \\
\\ 
\texttt{However, I also made a mistake by not considering the other player's perspective and the fact that we had already agreed on a deal. I should have proposed a secret proposal that takes into account the items we had already agreed on and also includes the crayon which is the most valuable item to both of us. } \\
\\ 
\texttt{A better approach would be to make a secret proposal that includes the items we had already agreed on and also includes the crayon, but also takes into account the other player's perspective and the fact that we had already agreed on a deal. } \\
\\ 
\texttt{To achieve this, I will make sure to carefully review the previous conversation and the items we had already agreed on. I will then propose a secret proposal that includes these items and also includes the crayon, while also considering the other player's perspective and the fact that we had already agreed on a deal. For example, I could have proposed a secret proposal like "{[}Proposal: 1 egg, 1 button, 1 crayon{]}" which includes the items we had already agreed on, and also includes the crayon which is the most valuable item to both of us.} \\
            }
        }
    }
    & & \\ \\

    \theutterance \stepcounter{utterance}  
    & & & \multicolumn{4}{p{0.6\linewidth}}{
        \cellcolor[rgb]{0.95,0.95,0.95}{
            \makecell[{{p{\linewidth}}}]{
                \texttt{\tiny{[P1$\rightarrow$P2]}}
                \texttt{Post{-}episode Reflexion (not sent to GM):} \\
\texttt{New plan: } \\
\\ 
\texttt{In my last attempt, I made a few mistakes that led to the game being aborted. Firstly, I didn't keep track of the items that were already proposed and agreed upon, which led to the other player asking for the second egg in exchange for the second button, even though I had already agreed to give them the other egg in exchange for the second button. } \\
\\ 
\texttt{To avoid this mistake, next time I will keep a mental or written record of the items that have been proposed and agreed upon, and make sure to refer to this record when making subsequent proposals or counter{-}proposals.} \\
\\ 
\texttt{Secondly, I didn't use the square bracket format correctly when making a secret proposal. } \\
\\ 
\texttt{To avoid this mistake, next time I will make sure to use the square bracket format correctly when making a secret proposal, and avoid using it when communicating with the other player.} \\
\\ 
\texttt{Lastly, I didn't think carefully enough about the implications of my proposals and counter{-}proposals, which led to the game being aborted.} \\
\\ 
\texttt{To avoid this mistake, next time I will take more time to think carefully about the implications of my proposals and counter{-}proposals, and make sure to consider the other player's perspective and the rules of the game.} \\
\\ 
\texttt{However, I also realized that I didn't make a secret proposal at all, which is a crucial part of the game. } \\
\\ 
\texttt{To avoid this mistake, next time I will make a secret proposal at the end of the game, after all the negotiations have taken place, and make sure to use the square bracket format correctly.} \\
\\ 
\texttt{Additionally, I should have focused more on finding a mutually beneficial agreement that maximizes the sum of our scores, rather than just trying to get the items I value most.} \\
\\ 
\texttt{To avoid this mistake, next time I will try to find a compromise that benefits both players, and make sure to consider the other player's perspective and the rules of the game.} \\
\\ 
\texttt{Here is a revised plan for the next attempt:} \\
\\ 
\texttt{1. Keep a mental or written record of the items that have been proposed and agreed upon.} \\
\texttt{2. Use the square bracket format correctly when making a secret proposal.} \\
\texttt{3. Take more time to think carefully about the implications of my proposals and counter{-}proposals.} \\
\texttt{4. Consider the other player's perspective and the rules of the game.} \\
\texttt{5. Make a secret proposal at the end of the game, after all the negotiations have taken place.} \\
\texttt{6. Focus on finding a mutually beneficial agreement that maximizes the sum of our scores.} \\
\\ 
\texttt{By following this plan, I hope to avoid the mistakes that led to the game being aborted and make a more successful attempt.} \\
            }
        }
    }
     \\ \\

\end{supertabular}
}

\end{document}
